\section{Способы представления графа}
\label{sec:approaches}

\subsection{Рамки решения}
\label{sec:scope}
Прежде чем выбирать представление, стоит зафиксировать рамки, в которых решается задача. Рёбер
$\bigO(m)$ --- это и есть <<гигабайт>> из условия, они не обязаны помещаться в память и живут на
диске. А структуры размером с число вершин, $\bigO(n)$ (векторы значений, степени, указатели),
наоборот, обязаны помещаться в ОЗУ: без них повершинную метрику не посчитать. Граница такая ---
$\bigO(m)$ не ограничено, $\bigO(n)$ обязано влезать.

Отсюда сразу отпадают подходы, которые выгружают на диск даже сами векторы значений $\bigO(n)$,
например Parallel Sliding Windows из GraphChi. Для этой задачи они избыточны: при <<гигабайте>> рёбер
вершин порядка $2$--$3$ млн, а их векторы --- это лишь десятки МБ, и они спокойно помещаются в память.
Случай, когда не влезает даже $\bigO(n)$, --- это уже другая задача, и она вне рамок данной работы.

\subsection{Push или pull}
Ядро метрики --- сумма по рёбрам~\eqref{eq:pull}. Организовать этот проход можно двумя способами, и
именно способ обработки определяет, как выгодно хранить граф.

\paragraph{Push (по источнику).} Обход идёт по рёбрам, сгруппированным по источнику $j$, и
$\mathrm{contrib}[j]$ прибавляется в ячейку приёмника $s_{\mathrm{new}}[i]$. Хранение здесь простое:
рёбра можно держать как есть, в порядке поступления, сортировать ничего не надо. Но при параллельном
проходе несколько потоков добавляют в одну и ту же ячейку популярного приёмника одновременно ---
возникает гонка. Пришлось бы ставить блокировки или атомарные операции, а порядок атомарных сложений
не определён, поэтому результат менялся бы в младших разрядах от запуска к запуску.

\paragraph{Pull (по приёмнику).} Обход идёт по приёмникам $i$, и для каждого собирается сумма
$s_{\mathrm{new}}[i] = \mathrm{base} + \sum_{j\to i}\mathrm{contrib}[j]$. У каждой ячейки результата
ровно один писатель, гонок нет. Зато появляется цена на хранении: чтобы входящие рёбра приёмника $i$
лежали подряд, рёбра нужно один раз сгруппировать и отсортировать по приёмнику. Это разовая работа на
этапе предобработки.

\subsection{Выбор}
Выбран pull. Решающим стало требование воспроизводимости. Push не требует сортировки, но его гонки и
неопределённый порядок сложения несовместимы с условием <<результат не должен зависеть от
случайности>>. Pull стоит разовой сортировки по приёмнику, зато на каждой из десятков итераций даёт
проход без блокировок, последовательное чтение с диска и детерминированный результат. Разовая
сортировка окупается уже на первых итерациях.

Из выбора pull прямо следует формат хранения: граф группируется по приёмнику и сортируется. Как он
раскладывается на диске и как строится за несколько внешних проходов --- в разделе~\ref{sec:arch}.
